\section{System Modeling and Methodology}
\subsection{Overall System Architecture and Cyber-Physical Modeling}
The foundation of addressing communication bottlenecks in Networked Control Systems (NCS) lies in accurately modeling the interplay 
between the physical plant dynamics and digital communication infrastructure.
\subsubsection{Networked Control System Architecture}
The proposed NCS architecture comprises a closed-loop cyber-physical pipeline encompassing sensors, a bandwidth-limited communication channel, a remote controller, and physical actuators. 
Unlike traditional tethered robots with dedicated point-to-point analog connections, the sensory and control data in this modern architecture must traverse a shared, bandwidth-limited communication medium.
\begin{enumerate}
    \item \textbf{Sensing and Edge Allocation Node}: High-precision encoders continuously monitor the physical states of the manipulator (e.g., joint angles and angular velocities). Instead of transmitting this full-precision data periodically, an edge computing node actsas the "JCC Allocator". It evaluates the necessity of transmission (Temporal ETC) and dynamically assigns limited quantization bits to different state variables (Spatial Allocation).
    \item \textbf{Bandwidth-Limited Communication Channel}: The network medium is characterized by strict capacity constraints, modeled as a maximum allowable payload of $B_{total}$ bits per sampling interval. This constraint applies universally, representing physical limits in transmission or artificial rate limits in shared wired networks.
    \item \textbf{Controller and Zero-Order Hold (ZOH)}: The remote controller receives the heavily quantized or delayed state information. A Zero-Order Hold (ZOH) mechanism is implemented at the receiver side to maintain the control input using the most recently received packet when new data is unavailable (due to ETC silence or network loss).
    \item \textbf{Actuators}: The calculated control torques are transmitted back to robot's joints motors to execute the desired trajectory.
\end{enumerate}
%% TODO: 加入图表
%% 这是一个大图 [图 3.1: 联合通信与控制系统 (JCC) 总体信息流架构图] (建议使用全页宽度的流程图，展示数据从真实传感器采集、边缘计算节点执行 ETC 判定与压缩、无线传输到接收端解压、控制电机输出的全过程)。
\subsubsection{Robotic Manipulator Dynamics and State-Space Modeling}
To derive control-oridented communication strategies, a rigorous mathematical model of the physical plant is required.
The continuous-time rigid-body dynamics of an $n$-degree-of-freedom (DOF) robotic manipulator can be described by the following 
nonlinear second-order differential equation:
\begin{equation}
    M(q)\ddot{q} + C(q,\dot{q})\dot{q} + G(q) = \tau
\end{equation}
where:
\begin{itemize}
    \item $q, \dot{q}, \ddot{q} \in \mathbb{R}^n$ represent the vectors of joint angles, angular velocities, and angular accelerations, respectively.
    \item $M(q) \in \mathbb{R}^{n \times n}$ is the symmetric, positive-definite inertial matrix.
    \item $C(q, \dot{q}) \in \mathbb{R}^{n \times n}$ represents the  Coriolis and centrifugal force matrix.
    \item $G(q) \in \mathbb{R}^n$ is the gravity loading vector.
    \item $\tau \in \mathbb{R}^n$ is the vector of control torques applied by the actuators.
\end{itemize}
To transform this into a standard state-space representation, let the system state vector be defined as $x(t) = [q(t)^T, \dot{q}(t)^T]^T \in \mathbb{R}^{2n}$,
and the control input as $u(t)=\tau(t)\in\mathbb{R}^n$. The nonlinear dynamics can be rewritten as $u(t) = \tau(t) \in \mathbb{R}^n$.
The nonlinear dynamics can be rewritten as a first-order ordinary differential equation $\dot{x}(t) = f(x(t), u(t))$:
\begin{equation}
    \dot{x}(t) = \frac{d}{dt} \begin{bmatrix} q \\ \dot{q} \end{bmatrix} = \begin{bmatrix} \dot{q} \\ M(q)^{-1}(\tau - C(q,\dot{q})\dot{q} - G(q)) \end{bmatrix}
\end{equation}
For advanced linear control methodologies (such as LQR) and sensitivity analysis, this nonlinear model is linearized around a nominal 
operating trajectory or equilibrium point $(x_e,u_e)$. Using Jacobian linearization, the continuous-time linear time-invariant (LTI) model is obtained:
\begin{equation}
    \dot{\tilde{x}}(t) = A_c \tilde{x}(t) + B_c \tilde{u}(t)
\end{equation}
where $x_k \in \mathbb{R}^{2n}$ is the discrete state at time step $k$,$u_k \in \mathbb{R}^n$ is the discrete control input, and 
$w_k \in \mathbb{R}^{2n}$ represents the additive process noise (e.g., modeling inaccuracies and physical disturbances), assumed to be 
zero-mean Gaussian noise with covariance $W$. The discrete matrices are $A = e^{A_c T_s}$ and $B = \int_{0}^{T_s} e^{A_c t} B_c dt$.
Physically, matrix $A$ captures the system's inertia and internal state couping, while matrix $B$ maps the applied actuator torques 
to the resulting state transitions.
\subsubsection{Bandwidth-Limited Communication Network Model}
The critical constraint in this thesis is the extreme limitation of the communication channel. Let the network capacity limit be 
strictly defined as a maximum allowed budget of $B_{total}$ bits per sampling interval $T_s$.

At any discrete time step $k$, if the sensor decides to transmit data, it can not send the ideal 64-bit floating-point representation of 
$x_k$. Instead, the state vector must be quantized. Let $b_(i,k)$ denote the number of bits allocated to the $i$-th element of the 
state vector $x_k$. The resource constraint mandates that:
\begin{equation}
    \sum_{i=1}^{2n} b_{i,k} \le B_{total}
\end{equation}
Let $\mathcal{Q}(\cdot)$ denote the quantization function. The transmitted and subsequently received quantized state 
is denoted as $x_k^q = \mathcal{Q}(x_k, b_k)$, introducing a quantization error $e_(q,k)=x_k-x_k^q$. If $b_(i,k)=0$, the 
$i$-th state is entirely discarded during transmission.

Furthermore, to model the temporal aspect the transmission, an indicator variable $\gamma_k\in\{0,1\}$ is introduced,
where $\gamma_k=1$ indicates a packet is triggered and successfully transmitted, and $\gamma_k=0$ indicates no transmission (event-triggered silence).

At the controller side, a Zero-Order Hold (ZOH) updates the estimated  state $\hat{x_k}$ based on the communication status:
\begin{equation}
    \hat{x}_k = 
    \begin{cases} 
    x_k^q, & \text{if } \gamma_k = 1 \text{ (Packet received)} \\
    \hat{x}_{k-1}, & \text{if } \gamma_k = 0 \text{ (No transmission, hold previous state)} 
    \end{cases}
\end{equation}
This can be concisely written as the receiver update equation:
\begin{equation}
    \hat{x}_k = \gamma_k x_k^q + (1 - \gamma_k) \hat{x}_{k-1}
\end{equation}
The control input $u_k$ is then strictly computed based on this remote estimate $\hat{x_k}$ rather the true physical state $x_k$.
The central challenge addressed in the following sections is how to simultaneously optimze $\gamma_k$ (when to transmit) and 
$b_k$ (how to allocate bits) to minimize the deviation between the actual physical trajectory and the desired reference.
\subsection{Temporal Domain: Event-Triggered Control Mechanism}
In Networked Control Systems, the temporal strategy of when to transmit data is as critical as the spatial strategy of how to 
compress it. This section details the deployment of an Event-Triggered Control (ETC) mechanism designed to aggressively truncate 
redundant transmissions in the temporal domain.
\subsubsection{The Inefficiency of Traditional Time-Triggered Sampling}
Traditionally, digital control systems rely on Time-Triggered Control (TTC), where sensor measurements are sampled and transmitted 
at rigidly fixed, periodic intervals $t_k=k T_s$. While TTC simplifies system design and analysis, it is notoriously inefficient 
regarding bandwidth utilization.

As rigorously demonstrated by Åström and Bernhardsson~\cite{method_crl}, Lebesgue sampling (event-based) significantly outperforms Riemann sampling (time-based) in stochastic control systems.
Under a TTC paradigm, the sensor continuously broadcasts data even when the robotic manipulator has reached a stable 
equilibrium or is executing a steady-state motion with negligible tracking error. In a strictly bandwidth-constrained 
shared network, shared network, this continuous stream of redundant, non-informative data squanders the limited channel capacity, 
directly exacerbating network congestion and packet loss for other critical nodes. To achieve "pragmatic communication," the system must abandon periodic transmission and adopt a paradigm of "speaking only when necessary".
\subsubsection{Design and Derivation of the Event-Triggering Condition}
The fundamental logic of Event-Triggered Control is to continuously evaluate the "novelty" or "urgency" of the current physical 
state at the sensor side and initiate a transmission only when the state deviation threatens the system's stability.

Let $\hat{x_k}$ represent the estimated state maintained at the remote controller via the Zero-Order Hold (ZOH). The measurement 
error (or event error), denoted as $e_k$, is defined as the discrepancy between the actual current state $x_k$ at the sensor 
and the remote estimate $\hat{x_k}$:
\begin{equation}
    e_k = x_k - \hat{x}_k
\end{equation}
The triggering condition is governed by a threshold mechanism. A transmission is triggered if and only if the norm of the 
measurement error exceeds a predefined safety threshold $\delta$:
\begin{equation}
    ||e_k||_2 > \delta
\end{equation}
\textbf{Mathematical Derivation of Stability}: The viability of this threshold is not arbitrary; it is strictly bounded by Lyapunov stability theory.
Based on the ETC framework established by Tabuada~\cite{method_etrs}, we must ensure that the control system maintains 
Asymptotic Stability despite the discontinuous state updates.

Assume there exits a stabliziing state-feedback controller $u_k=-K \hat{x_k}$, and a corresponding discrete-time Control Lyapunov 
Function $V(x_k)=x_k^T P_L x_k$ (where $P_L$ is a symmetric positive-define matrix). For the closed-loop system to be asymptotically
stable, the Lyapunov difference must be strictly negative define:
\begin{equation}
    \Delta V = V(x_{k+1}) - V(x_k) < 0
\end{equation}
By substituting the closed-loop dynamics $x_{k+1}=Ax_k-BK\hat{x_k}=(A-BK)x_k-BKe_k$ into the Lyapunov difference, $\Delta V$ can be 
bounded by a function of the true state $x_k$ and the measurement error $e_k$. According to Input-to-State Stability (ISS) properties,
$\Delta V$ takes the form:
\begin{equation}
    \Delta V \le -\alpha(||x_k||) + \beta(||e_k||)
\end{equation}
where $\alpha(\cdot)$ and $\beta(\cdot)$ are class $\mathcal{K}$ functions. 
To guarantee $\Delta V < 0$, the error must satisfy $\beta(||e_k||) < \alpha(||x_k||)$. 
This implies that as long as the measurement error $e_k$ is bounded by a relative threshold (e.g., $||e_k|| \le \sigma ||x_k||$) or a sufficiently small absolute static threshold $\delta$, the stability of the robotic manipulator is mathematically guaranteed.

In our proposed framework, to simplify the edge computation load on the sensor node, a robust absolute static threshold $||e_k||_2>\delta$
is implemented as the primary transmission trigger.
\subsubsection{Mathematical Abstraction of Communication Behavior}
To formally integrate the temporal ETC behavior with the spatial quantization algorithm (discussed in Section 3.3), we introduce 
a Boolean communication decision variable $\gamma_k\in\{0,1\}$.

The edge sensor node continuously evaluates the ETC condition and generates the decision variable:
\begin{equation}
    \gamma_k = 
    \begin{cases} 
    1, & \text{if } ||x_k - \hat{x}_{k-1}||_2 > \delta \quad \text{(Trigger Event)} \\
    0, & \text{otherwise} \quad \text{(Silence)}
    \end{cases}
\end{equation}
Consequently, the mathematical model of the receiver (the Zero-Order Hold at the controller) is perfectly unified. Upon integrating 
the quantization process $x_k^q = \mathcal{Q}(x_k, b_k)$, the remote state estimate $\hat{x_k}$ is updated according to the following 
equation.
\begin{equation}
    \hat{x}_k = \gamma_k x_k^q + (1 - \gamma_k) \hat{x}_{k-1}
\end{equation}
This equation powerfully summarizes the JCC temporal-spatial interplay: If $\gamma_k=1$, the network is actively utilized,
and the spatial quantizer $\mathcal{Q}(\cdot)$ determines the precision of the payload $x_k^q$. If $\gamma_k=0$, exactly zero bits
are transmitted over the channel, and the controller relies purely on the historical memory $\hat{x}_{k-1}$, thereby achieving massive bandwidth savings in the temporal domain.
\subsubsection{Spatial Domain: Static LQR-Weighted Data Compression}
While the temporal ETC mechanism dictates when to transmit, the spatial compression mechanism dictates \textit{how} to
allocate the severely limited bit budget $B_{total}$ across multiple joint states during an active transmission. Traditional 
approaches uniformly divide the bit budget ($b_i = B_{total} / n$), which blindly treats a 0.1-rad error at the base joint as equivalent
to a 0.1-rad error at the end-effector.

To overcome this, a static control-aware baseline utilizes the infinite-horizon Linear Quadratic Regulator (LQR) cost function
$J = \sum (x_k^T Q x_k + u_k^T R u_k)$. By solving the Discrete Algebraic Riccati Equation (DARE), the system extracts the Hessian 
sensitivity matrix $P$. According to rate-distortion theory, the optimal bit allocation $b_i$ that minimizes the control-oriented 
distortion is formulated as:
\begin{equation}
    b_i = \frac{B_{total}}{n} + \frac{1}{2} \log_2 \left( \frac{P_{ii} \sigma_i^2}{\left( \prod_{j=1}^n P_{jj} \sigma_j^2 \right)^{1/n}} \right)
\end{equation}
This static LQR-weighted allocation significantly improves control resilience by permanently stripping bits from low-sensitivity joints and granting them to high-sensitivity joints (e.g., base links). 
However, in highly nonlinear robotic dynamics, the true sensitivity of a joint is state-dependent and fluctuates drastically based on the manipulator's real-time posture and velocity. 
To achieve true context-awareness, the system must transition from static mathematical allocation to dynamic, learning-based arbitration.
\subsection{Deep Reinforcement Learning for Pragmatic Communication}
The core innovation of this thesis is the formulation of a dynamic, context-aware bandwidth allocator driven by Deep Reinforcement Learning (DRL). 
Rather than relying on fixed linear approximations (such as the static $P$ matrix), an AI agent is deployed at the edge sensor node to act as a centralized communication arbiter. 
The agent observes the real-time physical tracking errors and dynamically redistributes the strictly limited $B_{total}$ bits at every time step. 
This learning-based paradigm epitomizes "Pragmatic Communication"—the communication protocol evolves explicitly to serve the downstream physical task, learning to proactively sacrifice the precision of robust linkages to guarantee the survival of critical joints.
\subsubsection{Markov Decision Process (MDP) Formulation}
To enable DRL training, the dynamic bandwidth allocation problem must be strictly formulated as a Markov Decision Process (MDP), 
defined by the tuple $\langle \mathcal{S}, \mathcal{A}, \mathcal{P}, \mathcal{R}, \gamma \rangle$. \\\\
\textbf{1. State Space (Observation) $\mathcal{S}$}:
For the agent to make context-aware decisions, the observation space must encompass sufficient information regarding 
the current physical urgency of the robotic manipulator. Relying solely on the raw state $x_t$ is insufficient, as the agent 
would remain ignorant of the control objective. Therefore, the observation $s_t\in\mathcal{S}$ is constructed as an 8-dimensional 
continuous vector containing the real-time physical states, the target reference states, and the current tracking errors:
% TODO:我的花体是不是太"花"了
\begin{equation}
    s_t = \Big[ q_t^T, \ \dot{q}_t^T, \ q_{target, t}^T, \ (q_t - q_{target, t})^T \Big]^T \in \mathbb{R}^8
\end{equation}
This augmented state space allows the neural network to implicitly infer both the current nonlinear posture of the 
manipulator and the instantaneous tracking deviation, which fundamentally dictate the real-time sensitivity of each joint.\\\\
\textbf{2. Action Space $\mathcal{A}$ and Discrete Bit Mapping}:
The primary objective of the agent is to allocate the integer bit budget $B_{total}$ across the $n$ state variables.
Directly outputting discrete integers via reinforcement learning introduces non-differentiable operations and combinatorial 
explosion, severely hindering convergence. To resolve this, a continuous action space is adopted, combined with a deterministic 
Softmax mapping layer.

The DRL agent outputs a continuous action vector $a_t = [a_{1,t}, a_{2,t}, \dots, a_{n,t}]^T \in [-1, 1]^n$. To map these 
continuous unbounded preferences into a strict constraint $sum_{i=1}^n b_i = B_{total}$ where $b_i \in \mathbb{Z}^+$, 
the following transformation is applied:

First, the actions are converted into a probability distribution $p_i$ using Softmax function:
\begin{equation}
    p_i = \frac{\exp(a_{i,t})}{\sum_{j=1}^n \exp(a_{j,t})}
\end{equation}
Next, the probability distribution is scaled by the total bit budget to obtain the floating-point bit allocation:
\begin{equation}
    b_{float, i} = p_i \times B_{total}
\end{equation}
Finally, the floating-point values are rounded to the nearest integer $b_i=[b_{float},i]$. To strictly enforce the resource 
constraint due to rounding sequentially assigned to the state variables exhibiting the largest fractional remainder $(b_{float, i} - b_i)$;
conversely, if the sum is excessive, bits are iteratively deducted from the states holding the maximum allocation. This mapping 
guarantees that the strict network constraint is never violated while maintaining a smooth gradient landscape for DRL training.\\\\
\textbf{3. Reward Function Design $\mathcal{R}$}:
The reward function is the guiding signal that shapes the "pragmatic" intelligence of the agent. Since the 
bandwidth constraint $B_{total}$ is strictly enforced by the action mapping layer, there is no need to penalize the agent 
for bandwidth consumption during spatial allocation. Instead, the reward $r_t$ is purely and directly tied to the downstream 
physical control performance.

At each time step, the instantaneous reward is defined as the negative weighted tracking error (incorporating the LQR $Q$ matrix):
\begin{equation}
    r_t = - \left( (x_t - x_{target, t})^T Q (x_t - x_{target, t}) \right)
\end{equation}
For a 2-DOF manipulator where base joint is exponentially more critical than the shoulder joint, the $Q$ matrix 
is highly skewed (e.g., $Q=\text{diag}(1000,1)$). Because the DRL seeks to maximize the cumulative expected reward 
$\mathbb{E} [\sum \gamma^t r_t]$, it quickly discovers a counter-intuitive but mathematically optimal strategy: it will
willingly output actions that allocate 0 bits to the shoulder joint (allowing its local error to accumulate), in order 
to secure maximum bits for the base joint to prevent the massive--1000 penalty multiplier. This mechanism is the mathematical 
engine behind "Pragmatic Communication". 
\subsubsection{PPO Algorithm and Policy Network Architecture}
Given the continuous nature of the action space and requirement for stable monotonic convergence, the Proximal Policy Optimization (PPO)
algorithm~\cite{method_ppoa} is selected over value-based methods (like DQN) or traditional Policy Gradients (like DDPG). PPO 
operates under an Actor-Critic architecture, utilizing two separate neural networks:
\begin{enumerate}
    \item \textbf{The Actor Network (Policy Function $\pi_\theta(a_t | s_t)$)}: Parameterized by weights $\theta$, this network maps the current state $s_t$ to Gaussian probability distribution over the continuous actions $a_t$, from which the agent samples its allocation decisions.
    \item \textbf{The Critic Network (Value Function $V_\phi(s_t)$)}: Parameterized by weights $\phi$, this network estimates the expected cumulative discounted reward from state $s_t$, serving as baseline to evaluate the quality of the Actor's actions.
\end{enumerate}
Both the Actor and Critic networks are constructed using Multilayer Perceptrons (MLPs). The input layer ingests the 8-dimensional 
state vector $s_t$. This is followed by two hidden layers, each comprising 128 neurons activated by the ReLU (Rectified Linear Unit) 
non-linearity, ensuring sufficient representational capacity to capture the highly nonlinear relationship between robotic kinematics 
and bandwidth sensitivity.
% TODO: DRL图
To update the Actor network, PPO defines a clipped surrogate objective function. Left the advantage function $\hat{A}_t = R_t - V_\phi(s_t)$
represent how much better an action was compared to Critic's expectation. Let $r_t(\theta) = \frac{\pi_\theta(a_t | s_t)}{\pi_{\theta_{old}}(a_t | s_t)}$
denote the probability ratio between the updated policy and the old policy. The core objective maximized by PPO is:
\begin{equation}
    L^{CLIP}(\theta) = \hat{\mathbb{E}}_t \left[ \min \left( r_t(\theta)\hat{A}_t, \ \text{clip}(r_t(\theta), 1-\epsilon, 1+\epsilon)\hat{A}_t \right) \right]
\end{equation}
where $\epsilon$ is a hyperparameter (typically 0.2). The $\text{clip}$ function restricts the policy update step size.
If a new allocation strategy yields exceptionally high rewards, PPO will not greedily shift the neural network weights 
too drastically, preventing the policy from collapsing into suboptimal deterministic extremes (e.g., permanently assigning all bits to one joint regardless of the state).
This trust-region optimization is paramount for training stability in cyber-physical environments, where extreme exploration can lead to immediate system divergence.
\subsubsection{Training Stability and Hard Negative Mining Strategy}
Training a DRL agent in a closed-loop control system presents severe challenges regarding state exploration. If the initial 
DRL agent randomly allocates bits uniformly, the LQR controller might successfully stabilize the robot near the target trajectory.
In this safe, steady-state region, the errors are negligible, and the rewards $r_t$ approach zero. The agent suffers from a "vanishing signal problem"--
it learns that \textit{any} bit allocation works fine when robot is safe, and thus fails to learn the critical "survival prioritization"
required during extreme physical disturbances.

To overcome this, a \textbf{Hard Negative Mining (Edge State Initialization)} strategy is explicitly implemented 
within the training pipeline. During the \inlinecode{reset()} phase of each training episode, the robotic manipulator is not initialized at the safe origin or along the reference trajectory. 
Instead, the joints are deliberately spawned at extreme adversarial positions (e.g., $q_1, q_2$ injected with maximum permissible initial deviations).

This adversarial initialization forces the environment to immediately yield catastrophic physical costs (massive negative rewards) 
at $t=0$. To survive the episode without triggering the safety truncation threshold (which aborts the episode and administers an additional--500 penalty), 
the PPO agent is forced to rapidly discover the causal link between base-joint tracking errors and system collapse. It effectively 
"learns under duress", heavily accelerating the convergence towards the optimal pragmatic communication strategy.
